\section{An elementary covering lemma}
\label{sec:covering}

We seek an inexpensive cover of all sets containing many edges of a
uniform family with bounded total cost. The construction first selects
subsets lying in many edges. The remaining edges then have controlled
extension cost, while pair sampling makes long tuples unlikely to survive
uncovered: they must avoid quadratically many independent choices.
Unions of sampled edges also
appear in the covering constructions of Fischer and Person
\cite[Sections~1.1 and~2]{FischerPerson25}; the argument below is self-contained.

\begin{lemma}\label{lem:covering}
Let \(X\) be a finite set, let \(0<\rho<1\), and let \(r,h\) be integers
with \(1\le r\le h\) and \(h\ge64\). Suppose that
\(\mathcal H\subseteq\binom{X}{r}\) satisfies
\[
|\mathcal H|\rho^r\le2^{h+1}.
\]
Then the family
\[
\left\{S\subseteq X:
  |\{e\in\mathcal H:e\subseteq S\}|\ge2^{12h}\right\}
\]
has a cover \(\mathcal G\subseteq2^X\) with
\[
c_\rho(\mathcal G)\le3\cdot2^{-h}.
\]
\end{lemma}

\begin{proof}
Put
\[
m=2^{12h},\qquad T=2^{4h},\qquad \pi=2^{-8h}.
\]
For \(I\subseteq X\) with \(|I|\le r\), define its weighted degree by
\[
d(I)=|\{e\in\mathcal H:I\subseteq e\}|\rho^{r-|I|}.
\]

\emph{Large weighted degrees.}
Begin with the generators
\[
\mathcal C_0=\{I\subseteq X:|I|\le r,\ d(I)>T\}.
\]
Counting the subsets of each edge gives
\[
\sum_{\substack{I\subseteq X\\|I|\le r}}d(I)\rho^{|I|}
=2^r|\mathcal H|\rho^r.
\]
Consequently,
\begin{equation}\label{eq:covering-initial-cost}
c_\rho(\mathcal C_0)
\le T^{-1}2^r|\mathcal H|\rho^r
\le2^{r+1-3h}\le2^{1-2h}.
\end{equation}
Call an edge \emph{regular} if it contains no member of \(\mathcal C_0\),
and let \(\mathcal H'\) be the family of regular edges. In particular,
\[
I\subseteq e\in\mathcal H'
\quad\Longrightarrow\quad d(I)\le T.
\]

We will use this regularity in the following form. Given \(B\subseteq X\),
group the edges \(e\in\mathcal H'\) according to \(I=e\cap B\).
Each nonempty group contributes
\[
\sum_{\substack{e\in\mathcal H'\\e\cap B=I}}
  \rho^{|e\setminus B|}
=|\{e\in\mathcal H':e\cap B=I\}|\rho^{r-|I|}
\le d(I)\le T.
\]
It follows that
\begin{equation}\label{eq:covering-extension}
\sum_{e\in\mathcal H'}\rho^{|e\setminus B|}
\le T\sum_{i=0}^r\binom{|B|}{i}.
\end{equation}

\emph{Sampled pairs.}
Independently for each ordered pair \((e,e')\in(\mathcal H')^2\),
including the diagonal pairs, sample that pair with probability \(\pi\).
Let \(\mathcal C_1\) consist of the unions of the sampled pairs. Applying
\eqref{eq:covering-extension} with \(B=e\), so that \(|B|=r\), yields
\[
\sum_{e'\in\mathcal H'}\rho^{|e'\setminus e|}\le2^rT.
\]
Thus
\begin{align}
\mathbb E[c_\rho(\mathcal C_1)]
&\le\pi\sum_{e,e'\in\mathcal H'}\rho^{|e\cup e'|}
\notag\\
&\le\pi|\mathcal H'|\rho^r2^rT
\le2^{r+1-3h}\le2^{1-2h}.
\label{eq:covering-pair-cost}
\end{align}
The first inequality counts unions with multiplicity. Distinct ordered
pairs remain independent sampling indices even when their unions coincide.

\emph{Residual generators.}
For every ordered tuple \((e_1,\ldots,e_m)\) of distinct regular edges,
add \(e_1\cup\cdots\cup e_m\) to \(\mathcal C_2\) if none of the pairs
\((e_i,e_j)\), \(1\le i,j\le m\), was sampled.
Then
\[
\mathcal G=\mathcal C_0\cup\mathcal C_1\cup\mathcal C_2
\]
covers the target family for every sampling outcome. Indeed, a target set
\(S\) containing no generator from \(\mathcal C_0\) contains at least \(m\)
distinct regular edges. Choose \(m\) of them. If a pair among these edges
was sampled, its union is a member of \(\mathcal C_1\) contained in \(S\);
otherwise their full union belongs to \(\mathcal C_2\).

It remains to bound the expected cost of \(\mathcal C_2\). For
\(0\le t\le m\), let
\[
W_t=
\sum_{\substack{(e_1,\ldots,e_t)\in(\mathcal H')^t\\
                e_1,\ldots,e_t\ \text{distinct}}}
  \rho^{|e_1\cup\cdots\cup e_t|},
\qquad W_0=1.
\]
For a prefix of length \(t<m\), its union \(B\) has size at most \(rt\le rm\).
The elementary inequality
\[
\sum_{i=0}^r\binom{|B|}{i}\le(1+|B|)^r\le(1+rm)^r
\]
can be seen by ordering \(B\), listing the members of each subset of size
at most \(r\), and padding the list to length \(r\) with a new symbol.
Since
\(\rho^{|B\cup e|}=\rho^{|B|}\rho^{|e\setminus B|}\),
\eqref{eq:covering-extension} bounds the total weight of extensions of
this prefix by
\[
\rho^{|B|}\,T(1+rm)^r.
\]
Allowing all regular edges as possible extensions only increases this
bound. Summing over prefixes and iterating gives
\[
W_{t+1}\le T(1+rm)^rW_t,
\qquad
W_m\le[T(1+rm)^r]^m.
\]

The \(m^2\) ordered pair indices associated with a tuple of distinct
edges are all distinct. Hence the probability that this tuple contributes
a residual generator is exactly
\[
(1-\pi)^{m^2}\le e^{-\pi m^2}\le2^{-\pi m^2}.
\]
Linearity of expectation, again counting generators with multiplicity,
therefore gives
\[
\mathbb E[c_\rho(\mathcal C_2)]
\le[T(1+rm)^r]^m\,2^{-\pi m^2}.
\]
To make this at most \(2^{-m}\), it is enough that
\[
\log_2 T+r\log_2(1+rm)+1\le\pi m.
\]
The choice of \(m\) makes \(\pi m=2^{4h}\), while the left side
is bounded by a quadratic function of \(h\). Indeed,
\[
1+rm\le(h+1)2^{12h}\le2^{13h},
\qquad
4h+13rh+1\le18h^2
\le18\cdot2^{2h}\le2^{4h}=\pi m.
\]
Here we used \(r\le h\), \(h+1\le2^h\), and \(h\ge64\).
It follows that
\begin{equation}\label{eq:covering-residual-cost}
\mathbb E[c_\rho(\mathcal C_2)]
\le2^{(4h+13rh)m-\pi m^2}\le2^{-m}.
\end{equation}

Combining \eqref{eq:covering-initial-cost},
\eqref{eq:covering-pair-cost}, and
\eqref{eq:covering-residual-cost}, we obtain
\[
\mathbb E[c_\rho(\mathcal G)]
\le2^{2-2h}+2^{-m}\le3\cdot2^{-h}.
\]
All families and sampling indices are finite, so some outcome has cost
at most its expectation. Since every outcome supplies a cover, this proves
the lemma.
\end{proof}
